\section{Extração de Características e Aprendizado de Máquina}
A classificação de gestos a partir de sinais sEMG requer a extração de características discriminativas que representem adequadamente os padrões musculares associados a cada movimento \cite{Castro}. No domínio do tempo, as características mais utilizadas incluem o valor RMS (Root Mean Square), o valor médio absoluto (MAV), a variância, a taxa de cruzamentos por zero (ZC) e o comprimento de forma de onda (WL). No domínio da frequência, a frequência mediana, a frequência média e a potência em sub-bandas específicas fornecem informações complementares sobre o recrutamento das unidades motoras.
Embora existam diversos algoritmos de classificação amplamente consolidados na literatura para sinais sEMG, como Máquinas de Vetores de Suporte (SVM), Redes Neurais Artificiais (RNA) do tipo MLP e modelos de Floresta Aleatória (Random Forest), a proposta deste trabalho consiste em realizar uma etapa prévia de filtragem dos sinais e, em seguida, utilizar Redes Neurais Convolucionais (CNNs), em consonância com a abordagem proposta em "Multiday EMG-Based Classification of Hand Motions with Deep Learning Techniques"  \cite{Rehman} . Essa estrutura apresenta a vantagem de remover ruídos e interferências na etapa de pré-processamento, permitindo que a CNN realize a extração automática de características diretamente das ondas filtradas.
